\documentclass{article}
\usepackage[T1]{fontenc}
\usepackage[utf8]{inputenc}
\usepackage{amsmath}
\usepackage{comment}
\usepackage{amssymb}
\usepackage{polski}
\usepackage[polish]{babel}
\usepackage{graphicx}
\usepackage{float} % Inf
\usepackage{enumitem}
\usepackage{subcaption}
\usepackage{yfonts}
% \usepackage{braket} % kwant
% \usepackage{tikz}
% \usetikzlibrary{positioning, arrows.meta, calc}
\usepackage[margin=2cm]{geometry}
% \usepackage{tabularx}
%\usepackage[version=4]{mhchem} % Chemia
\usepackage{hyperref}
\hypersetup{
colorlinks=true,
linkcolor=black,
urlcolor=red
}
\usepackage{listings}
\lstset{
  basicstyle=\ttfamily\small,
  breaklines=false,
  frame=single
}

\title{\Huge{\textgoth{Sprawozdanie - Projekt 7}}}
\author{\Huge{\textfrak{Maciej Flaga}}}
\date{}

\begin{document}
\maketitle

% \section{Wstęp teoretyczny}

\section{Realizacja}

\subsection{Metoda RS}

Zdefiniowano potencjał jako pojedynczą studnię i uwzględniono periodyczne warunki brzegowe w macierzy hamiltonianu. Napisano funkcję \texttt{diagHdlaK}, która korzysta z biblioteki \texttt{GSL} i jej funkcji do wyznaczania wartości własnych macierzy. Następnie zdiagonalizowano macierz dla różnych wartości $k$ na siatce $k\in[-\pi/a, \pi/a]$. Wykreślono dane poziomy energetyczne w funkcji wektora falowego:

\begin{figure}[H]
  \centering
  \includegraphics[width=0.85\linewidth]{mp.png}
  \caption{Wyniki dla $V_0=-100,\ a=1$}
\end{figure}

\section{Źródła}

\end{document}
